% 对偶空间：线性泛函的内积表示。

\section{线性泛函的表示}
\phantomsection
\label{sec:ls-riesz-representation}
在$\mathbb R^d$的欧氏内积下，加权读取$\varphi(\bm x)=\bm a^\top\bm x$就是与向量$\bm a$的内积。在一般有限维内积空间中，每个读取也都能这样写。
\begin{theorem}[有限维Riesz表示定理]
\label{thm:ls-riesz-representation}
设$V$是有限维内积空间。对每个$\varphi\in V^*$，存在唯一的$\bm v_\varphi\in V$，使$\varphi(\bm x)=\langle\bm x,\bm v_\varphi\rangle$对所有$\bm x\in V$成立。
\end{theorem}
\begin{proof}
取$V$的一组标准正交基$\bm q_1,\ldots,\bm q_d$（第~\ref{sec:ls-orthonormal-basis-construction}节），令$\bm v_\varphi=\sum_j\varphi(\bm q_j)\bm q_j$。由式~\eqref{eq:ls-orthonormal-expansion}，$\bm x=\sum_j\langle\bm x,\bm q_j\rangle\bm q_j$，于是
\[
 \begin{aligned}
 \varphi(\bm x)&=\sum_j\langle\bm x,\bm q_j\rangle\varphi(\bm q_j)\\
 &=\Bigl\langle\bm x,\sum_j\varphi(\bm q_j)\bm q_j\Bigr\rangle=\langle\bm x,\bm v_\varphi\rangle.
 \end{aligned}
\]
若$\bm v,\bm v'$都满足要求，则$\langle\bm x,\bm v-\bm v'\rangle=0$对一切$\bm x$成立，取$\bm x=\bm v-\bm v'$得$\bm v=\bm v'$。
\end{proof}
对应$\varphi\mapsto\bm v_\varphi$保持线性组合，这由表示向量的唯一性得到；每个$\bm v$又都表示读取$\langle\cdot,\bm v\rangle$。所以这个对应是$V^*$到$V$的双射。

在这个对应之下，标准正交基的对偶基就是这组基本身，因为第$j$个坐标读取是$\bm x\mapsto\langle\bm x,\bm q_j\rangle$。对一般的基，对偶基对应的向量与原基不同：图~\ref{fig:ls-dual-basis-readout}中的红色箭头$\bm b^1,\bm b^2$，正是对偶读取$\varphi_1,\varphi_2$在欧氏内积下对应的向量。

表示向量依赖内积，读取本身却不依赖。在$\mathbb R^d$上取加权内积$\langle\bm x,\bm y\rangle_{\bm M}=\bm x^\top\bm M\bm y$，同一个读取可以改写为
\[
 \varphi(\bm x)=\bm a^\top\bm x=\bm x^\top\bm M(\bm M^{-1}\bm a),
\]
所以它对应的向量是$\bm M^{-1}\bm a$，而不是$\bm a$。函数空间中也是如此。
\begin{example}[积分读取在两点内积下的表示向量]
\label{ex:ls-riesz-integral-two-point}
在$P_1$上取两点内积$\langle f,g\rangle=f(0)g(0)+f(1)g(1)$，考虑积分读取$\varphi(f)=\int_0^1f(t)\,\mathrm dt$。对$f=a+bt$，
\[
 \varphi(f)=a+\frac b2 .
\]

要把它写成与某个固定函数的内积，而这个内积只用到$f(0)=a$与$f(1)=a+b$，所以先把$a+b/2$表示成这两个值的组合：
\[
 a+\frac b2=\tfrac12\bigl(f(0)+f(1)\bigr).
\]
这正是$f$与常数函数$\tfrac12$的内积：
\[
 \Bigl\langle f,\tfrac12\Bigr\rangle=\tfrac12\bigl(f(0)+f(1)\bigr)=\tfrac12(a+a+b)=a+\frac b2 .
\]

因此$\bm v_\varphi$是常数函数$\tfrac12$。在系数向量的欧氏内积下，同一个读取对应$(1,1/2)^\top$，表示向量随内积而变。等式$\varphi(f)=\tfrac12\bigl(f(0)+f(1)\bigr)$也说明梯形公式对一次函数精确。
\end{example}

\paragraph{复空间中的表示}
复空间沿用前面“第一变量共轭线性、第二变量线性”的约定。这时线性泛函应写成$\varphi(\bm x)=\langle\bm v,\bm x\rangle$，让$\bm x$处于线性的变量位置。若$\bm q_1,\ldots,\bm q_d$是标准正交基，则
\[
 \bm v=\sum_j\overline{\varphi(\bm q_j)}\bm q_j ;
\]
代入各基向量即可核对读数。表示向量随泛函作共轭线性变化，这与实空间中的线性对应不同。

例如在$\mathbb C$的标准内积$\langle v,z\rangle=\overline v z$下，泛函$\varphi(z)=(1+\mathrm i)z$对应$v=1-\mathrm i$，因为$\overline v z=(1+\mathrm i)z$。若交换两个变量的位置，得到的表达式关于$z$共轭线性，不再是线性泛函。所以共轭来自内积约定，不能从实数公式中直接删去。
